\chapter{Частицы, поля и релятивистские обозначения}

\section{Задача о доплеровском сдвиге перед лекцией}

Обсуждение начинается с вопроса, оставшегося с предыдущей недели. Всегда ли наблюдатель, движущийся почти со скоростью света, измеряет большую длину волны? Нет: ответ зависит от направления движения. Наблюдатель, движущийся вместе со световой волной, видит, что ее длина волны увеличена, тогда как наблюдатель, движущийся навстречу волне, видит, что она сжата. Это релятивистский эффект Доплера.

\begin{classroomqa}
\audiencequestion{Увеличивает ли движение со скоростью, близкой к световой, наблюдаемую длину волны?}
\lecturerresponse{Только если наблюдатель движется в том же направлении, что и свет. Изменение знака скорости наблюдателя меняет знак сдвига. Этот результат можно получить непосредственно из геометрии Лоренца, не заучивая отдельную формулу. \lecturetimestamp{00:00:13}}
\end{classroomqa}

Используем единицы, в которых \(c=1\), и пусть штрихованная система отсчета движется вправо со скоростью \(v\). Гребень световой волны, движущийся вправо, следует светоподобной прямой. Нормируем длину волны в нештрихованной системе к единице и запишем мировые линии двух последовательных гребней как \begin{equation}
x=t,
\qquad
x=t+1.
\label{eq:sr5-crest-lines}
\end{equation} Длина волны в нештрихованной системе измеряется между гребнями при одинаковом \(t\). В штрихованной системе ее нужно измерять между ними при одинаковом \(t'\). Поэтому весь расчет сводится к относительности одновременности.

\begin{figure}[htbp]
\centering
\includegraphics[width=0.90\textwidth]{lecture_05_figure_01.png}
\caption{Последовательные гребни волны, движущейся вправо, световой конус и наклоненные оси движущейся системы отсчета.}
\lectureframe{00:03:30}
\end{figure}

На этом месте вывод ненадолго прерывается обсуждением последовательности курсов. Летнего курса не будет. И общая теория относительности, и более широкое возвращение к квантовой механике рассматриваются как кандидаты на следующий курс; при этом гармонический осциллятор выделяется как материал, который следует понять до квантовой теории поля.

\begin{classroomqa}
\audiencequestion{Мог бы один из будущих курсов обсудить инфракрасно--ультрафиолетовую связь?}
\lecturerresponse{Этот вопрос относится к стыку квантовой механики и гравитации. У слушателей уже достаточно квантовой механики, чтобы распознать проблему, но сначала нужно развить гравитацию, естественнее всего через общую теорию относительности. \lecturetimestamp{00:05:37}}
\end{classroomqa}

\begin{classroomqa}
\audiencequestion{Можно ли просто считать относительный масштаб наклоненных осей с рисунка на доске?}
\lecturerresponse{Нет. Пространственно-временная диаграмма отражает причинную структуру и геометрию одновременности, но обычная евклидова линейка на рисунке не является лоренцевым измерением. Нужно определить соответствующие события и применить преобразование Лоренца. \lecturetimestamp{00:06:10}}
\end{classroomqa}

Теперь расчет продолжается. Нормируем расстояние между гребнями в нештрихованной системе к единице и алгебраически найдем два одновременных в штрихованной системе события.

Преобразование Лоренца имеет вид \begin{equation}
t'=\gamma(t-vx),
\qquad
x'=\gamma(x-vt),
\qquad
\gamma=\frac{1}{\sqrt{1-v^2}}.
\label{eq:sr5-lorentz-doppler}
\end{equation} Выберем срез одновременности в штрихованной системе \(t'=0\). Тогда из уравнения \eqref{eq:sr5-lorentz-doppler} следует \begin{equation}
t=vx.
\label{eq:sr5-primed-simultaneity}
\end{equation} Первый гребень пересекает этот срез в начале координат. Второй гребень подчиняется \(x=t+1\), поэтому для его точки пересечения имеем \begin{align}
x&=vx+1,\\
x&=\frac{1}{1-v},
\qquad
t=\frac{v}{1-v}.
\end{align} Его штрихованная координата равна \begin{align}
x'
&=\frac{x-vt}{\sqrt{1-v^2}}\\
&=\frac{1-v^2}{(1-v)\sqrt{1-v^2}}\\
&=\sqrt{\frac{1+v}{1-v}}.
\end{align} Поскольку исходная длина волны была равна единице, \begin{equation}
\boxed{
\frac{\lambda'}{\lambda}
=
\sqrt{\frac{1+\beta}{1-\beta}}
},
\qquad
\beta=\frac{v}{c}.
\label{eq:sr5-doppler}
\end{equation} Эквивалентно, \begin{equation}
\frac{\lambda'}{\lambda}
=
\sqrt{\frac{c+v}{c-v}}.
\end{equation}

\begin{figure}[htbp]
\centering
\includegraphics[width=0.90\textwidth]{lecture_05_figure_02.png}
\caption{Доплеровский множитель после восстановления \(c\); изменение знака \(v\) дает обратный сдвиг.}
\lectureframe{00:13:40}
\end{figure}

При \(v>0\) наблюдатель в штрихованной системе догоняет свет и измеряет большую длину волны. Замена \(v\) на \(-v\) описывает движение навстречу свету и делает этот множитель меньшим единицы. Если бы наблюдатель, догоняя луч, мог приближаться к \(c\), длина волны расходилась бы. Сдвиг частоты дается обратным отношением.

\section{Почему взаимодействие должно входить в одно действие}

После предварительных вопросов тема — поля и частицы. В квантовой теории поля частицы и поля в конечном счете станут двумя описаниями одних и тех же возбуждений. Но здесь рамка не такова. Здесь частицы следуют по мировым линиям, а классические поля заполняют пространство-время. Ближайший вопрос скромнее: если поле действует на частицу, то как частица действует на поле в ответ?

Принцип действия отвечает на это на примере конечномерного прототипа. Пусть \(x(t)\) и \(y(t)\) — две обобщенные координаты. Если лагранжиан распадается, \begin{equation}
L(x,\dot x,y,\dot y)
=
L_x(x,\dot x)+L_y(y,\dot y),
\label{eq:sr5-separable-l}
\end{equation} то два уравнения Эйлера--Лагранжа независимы: \begin{equation}
\frac{d}{dt}\frac{\partial L}{\partial\dot x}
=
\frac{\partial L}{\partial x},
\qquad
\frac{d}{dt}\frac{\partial L}{\partial\dot y}
=
\frac{\partial L}{\partial y}.
\end{equation} Уравнение для \(x\) не содержит \(y\), а уравнение для \(y\) не содержит \(x\).

Чтобы заставить переменные взаимодействовать, добавим член, который нельзя разложить на \(x\)-часть и \(y\)-часть. Простой пример таков: \begin{equation}
L
=
\frac{\dot x^2}{2}
+\frac{\dot y^2}{2}
-V(x)-W(y)+\kappa xy.
\label{eq:sr5-coupled-coordinates}
\end{equation} Теперь варьирование \(x\) дает член, пропорциональный \(y\), а варьирование \(y\) дает член, пропорциональный \(x\). Один и тот же смешанный член управляет обоими направлениями влияния.

\begin{figure}[htbp]
\centering
\includegraphics[width=0.90\textwidth]{lecture_05_figure_03.png}
\caption{Две обобщенные координаты становятся динамически связанными, когда один член лагранжиана содержит их обе.}
\lectureframe{00:19:00}
\end{figure}

Именно такова точная формулировка утверждения о действии--противодействии, используемого здесь: оба объекта динамические, и оба варьируются в одном общем действии. Если рассматривать поле как внешне заданный фон, его отклик может скрыться, но как только само поле становится динамическим, член связи дает оба уравнения.

\section{Как скалярное поле действует на частицу}

\subsection{Релятивистская связь}

Сначала будем считать скалярное поле \(\phi(\mathbf{x},t)\) заданным. Это может быть волна или статический профиль; пока его собственное уравнение несущественно. Действие свободной релятивистской частицы имеет вид \begin{equation}
S_{\mathrm p}
=
-m\int d\tau.
\end{equation} Для движения по прямой и при \(c=1\) получаем \begin{equation}
L_{\mathrm p}^{(0)}
=
-m\sqrt{1-\dot x^2}.
\end{equation} Простейшая скалярная связь заменяет массовый множитель на зависящий от поля: \begin{equation}
\boxed{
L_{\mathrm p}
=
-\bigl[m+g\phi(x,t)\bigr]\sqrt{1-\dot x^2}
}.
\label{eq:sr5-particle-coupling}
\end{equation} Константа связи \(g\) измеряет, насколько сильно поле воздействует на эту частицу. Ее знак — дело соглашения: замена \(g\) на \(-g\) обращает направление скалярной силы.

\begin{figure}[htbp]
\centering
\includegraphics[width=0.90\textwidth]{lecture_05_figure_04.png}
\caption{Действие свободной частицы, в котором \(m\) заменено локальным коэффициентом \(m+g\phi(x,t)\).}
\lectureframe{00:24:10}
\end{figure}

\subsection{Предел малых скоростей частицы}

Восстановим скорость света, чтобы ньютоновский предел был виден явно: \begin{equation}
L_{\mathrm p}
=
-\bigl[mc^2+g\phi(x,t)\bigr]
\sqrt{1-\frac{\dot x^2}{c^2}}.
\label{eq:sr5-particle-coupling-c}
\end{equation} При \(|\dot x|\ll c\) \begin{equation}
\sqrt{1-\frac{\dot x^2}{c^2}}
=
1-\frac{\dot x^2}{2c^2}
+O\!\left(\frac{\dot x^4}{c^4}\right).
\end{equation} Подстановка в \eqref{eq:sr5-particle-coupling-c} дает \begin{align}
L_{\mathrm p}
&=
-mc^2-g\phi
+\frac{m\dot x^2}{2}
+\frac{g\phi\,\dot x^2}{2c^2}
+O(c^{-4})\\
&\longrightarrow
\frac{m\dot x^2}{2}-g\phi(x,t),
\label{eq:sr5-nr-particle}
\end{align} после отбрасывания постоянной \(-mc^2\) и членов, подавленных как \(c^{-2}\). Получается обычная форма \(T-V\), так что \begin{equation}
V_{\mathrm p}(x,t)=g\phi(x,t).
\label{eq:sr5-scalar-potential}
\end{equation}

\begin{figure}[htbp]
\centering
\includegraphics[width=0.90\textwidth]{lecture_05_figure_05.png}
\caption{Разложение корня оставляет ньютоновский кинетический член и скалярную потенциальную энергию \(g\phi\).}
\lectureframe{00:28:50}
\end{figure}

Большая постоянная энергии покоя не входит в уравнение Эйлера--Лагранжа. Зависящий от поля член входит. В электромагнетизме электрический заряд играет аналогичную роль параметра связи, хотя реальное электромагнитное взаимодействие использует четырехвекторный потенциал, а не эту скалярную игрушечную модель.

\section{Как частица действует на поле в ответ}

\subsection{Одна вариационная задача}

Теперь поле уже не может оставаться просто заданным. Рассмотрим область пространства-времени, содержащую и конфигурацию поля, и мировую линию частицы. Физическое решение делает полное действие стационарным по отношению к вариациям и поля, и мировой линии: \begin{equation}
\delta S[\phi,x]=0.
\end{equation} Варьирование мировой линии при фиксированном поле дает уравнение частицы. Варьирование поля при фиксированной мировой линии дает уравнение поля. Именно общий член взаимодействия и связывает их.

\begin{figure}[htbp]
\centering
\includegraphics[width=0.90\textwidth]{lecture_05_figure_06.png}
\caption{Мировая линия частицы, проходящая через область пространства-времени, занятую скалярным полем.}
\lectureframe{00:30:30}
\end{figure}

Для свободного безмассового скалярного поля возьмем \begin{equation}
S_{\mathrm f}
=
\int d^4x\,
\frac12
\left[
(\partial_t\phi)^2
-(\partial_x\phi)^2
-(\partial_y\phi)^2
-(\partial_z\phi)^2
\right].
\label{eq:sr5-free-field-action}
\end{equation} Начиная с этого места, снова используем \(c=1\). При явном сохранении единиц временной член нес бы соответствующий множитель \(1/c^2\), в зависимости от нормировки координат и поля.

\begin{figure}[htbp]
\centering
\includegraphics[width=0.90\textwidth]{lecture_05_figure_07.png}
\caption{Действие поля, записанное в той же области пространства-времени, что и мировая линия частицы.}
\lectureframe{00:32:10}
\end{figure}

\subsection{Неподвижная частица и интеграл по мировой линии}

Чтобы выявить источник с минимальной алгеброй, поместим покоящуюся частицу в начало координат. Тогда множитель \(d\tau\) равен \(dt\), и часть взаимодействия в \eqref{eq:sr5-particle-coupling} принимает вид \begin{equation}
S_{\mathrm{int}}
=
-g\int dt\,\phi(\mathbf{0},t).
\label{eq:sr5-worldline-interaction}
\end{equation} Это интеграл только вдоль мировой линии, тогда как \eqref{eq:sr5-free-field-action} — интеграл по всему пространству-времени. Чтобы объединить их в один пространственно-временной интеграл, нужна функция, выделяющая положение частицы.

\begin{figure}[htbp]
\centering
\includegraphics[width=0.90\textwidth]{lecture_05_figure_08.png}
\caption{Неподвижная мировая линия внутри области поля и действие поля, которое нужно с ней объединить.}
\lectureframe{00:35:35}
\end{figure}

\subsection{Дельта-функция Дирака}

Одномерная дельта-функция определяется соотношением
\begin{equation}
\int_{-\infty}^{\infty}dx\,f(x)\delta(x-a)=f(a).
\label{eq:sr5-delta-definition}
\end{equation}
Вне \(x=a\) она обращается в нуль, имеет особенность при \(x=a\) и обладает единичным интегралом.
Ее лучше понимать как распределение, характеризуемое \eqref{eq:sr5-delta-definition}, а не как обычную конечнозначную функцию.

\begin{figure}[htbp]
\centering
\includegraphics[width=0.90\textwidth]{lecture_05_figure_09.png}
\caption{Тождество с дельта-функцией, переписывающее значение поля в начале координат в виде пространственного интеграла.}
\lectureframe{00:36:45}
\end{figure}

В трех измерениях
\begin{equation}
\delta^3(\mathbf{x})
=
\delta(x)\delta(y)\delta(z),
\end{equation}
и, следовательно,
\begin{align}
\int d^3x\,
\phi(\mathbf{x},t)\delta^3(\mathbf{x})
&=
\int dx\,dy\,dz\,
\phi(x,y,z,t)\\
&\quad\times\delta(x)\delta(y)\delta(z)\\
&=\phi(\mathbf{0},t).
\label{eq:sr5-triple-delta}
\end{align}

\begin{classroomqa}
\audiencequestion{Откуда берется \(d^3x\)?}
\lecturerresponse{Это обычный элемент пространственного объема \(dx\,dy\,dz\). В одном пространственном измерении то же тождество содержит только \(dx\) и одну дельта-функцию. \lecturetimestamp{00:37:09}}
\end{classroomqa}

\begin{figure}[htbp]
\centering
\includegraphics[width=0.90\textwidth]{lecture_05_figure_10.png}
\caption{Трехмерная дельта-функция, представленная как произведение трех одномерных множителей.}
\lectureframe{00:39:05}
\end{figure}

Уравнение \eqref{eq:sr5-worldline-interaction} теперь можно записать в виде
\begin{equation}
S_{\mathrm{int}}
=
-g\int d^4x\,
\delta^3(\mathbf{x})\phi(\mathbf{x},t).
\end{equation}
Полное действие для статического источника имеет вид
\begin{align}
S
=
\int d^4x
\biggl\{
&\frac12(\partial_t\phi)^2
-\frac12|\boldsymbol{\nabla}\phi|^2
\nonumber\\
&-g\delta^3(\mathbf{x})\phi
\biggr\}.
\label{eq:sr5-combined-action}
\end{align}

\enlargethispage{1pt}
\begin{classroomqa}
\audiencequestion{Не потерялся ли знак минус, когда взаимодействие было подставлено в пространственно-временное действие?}
\lecturerresponse{Лагранжиан мировой линии содержит \(-g\phi\), поэтому пространственно-временная плотность должна содержать \(-g\delta^3(\mathbf{x})\phi\). Сохранение этого знака фиксирует член источника в уравнении поля. \lecturetimestamp{00:41:24}}
\end{classroomqa}

\begin{figure}[htbp]
\centering
\includegraphics[width=0.90\textwidth]{lecture_05_figure_11.png}
\caption{Плотность свободного скалярного поля и локализованное взаимодействие, объединенные в одно пространственно-временное действие.}
\lectureframe{00:42:05}
\end{figure}

\subsection{Уравнение поля с источником}

Варьирование \eqref{eq:sr5-combined-action} по \(\phi\) дает
\begin{equation}
\boxed{
\partial_t^2\phi-\nabla^2\phi
=
-g\delta^3(\mathbf{x})
}.
\label{eq:sr5-sourced-wave}
\end{equation}
Вне начала координат это свободное волновое уравнение. В начале координат дельта-функция указывает на частицу как на локализованный источник.

\begin{figure}[htbp]
\centering
\includegraphics[width=0.90\textwidth]{lecture_05_figure_12.png}
\caption{Волновой оператор с точечным источником в начале координат.}
\lectureframe{00:45:35}
\end{figure}

Для статического решения \(\partial_t\phi=0\), поэтому
\begin{equation}
\boxed{
\nabla^2\phi
=
g\delta^3(\mathbf{x})
}.
\label{eq:sr5-poisson}
\end{equation}
Это уравнение Пуассона для точечного источника. Аналогия здесь электростатическая, пока еще не электромагнитная: поле по-прежнему есть скаляр \(\phi\), а не электромагнитный четырехпотенциал.

\begin{figure}[htbp]
\centering
\includegraphics[width=0.90\textwidth]{lecture_05_figure_13.png}
\caption{Статический предел волнового уравнения с источником, распознаваемый как уравнение Пуассона.}
\lectureframe{00:47:05}
\end{figure}

Теперь взаимность выражена явно. Член \(g\phi\) придает частице потенциальную энергию, и то же \(g\) делает частицу источником для \(\phi\). Эти два эффекта не являются независимыми предположениями.

\subsection{Движущиеся источники излучают}

Частица, фиксированная в точке \(\mathbf{a}\), описывается \(\delta^3(\mathbf{x}-\mathbf{a})\). Если ее положение меняется со временем, источник принимает вид
\begin{equation}
\delta^3\!\left(\mathbf{x}-\mathbf{a}(t)\right).
\label{eq:sr5-moving-delta}
\end{equation}
Зависящая от времени правая часть не может уравновешиваться полем, не зависящим от времени. Член \(\partial_t^2\phi\) должен вернуться, и ускоряющийся или колеблющийся источник может порождать волны в поле.

\begin{figure}[htbp]
\centering
\includegraphics[width=0.90\textwidth]{lecture_05_figure_14.png}
\caption{Источник с дельта-функцией, смещенный в \(\mathbf{a}(t)\), что делает уравнение поля зависящим от времени.}
\lectureframe{00:50:30}
\end{figure}

\begin{classroomqa}
\audiencequestion{Если источник рассматривается нерелятивистски, то где в уравнении поля проявляется релятивистский характер?}
\lecturerresponse{Релятивистское распространение закодировано в пространственно-временном волновом операторе слева. Полностью релятивистское описание частицы сохраняло бы ее действие по собственному времени и скоростные множители. Даже медленно движущийся, но зависящий от времени источник может возбуждать волны. \lecturetimestamp{00:52:14}}
\end{classroomqa}

\begin{classroomqa}
\audiencequestion{Какую размерность должно иметь \(g\phi\)?}
\lecturerresponse{Оно стоит рядом с \(mc^2\) в действии частицы, поэтому имеет размерность энергии. В нормировке с естественными единицами, используемой здесь, скаляр в \(3+1\) измерениях имеет массовую размерность 1, и \(g\) может быть безразмерным. Размерности меняются с размерностью пространства-времени и с нормировкой поля. \lecturetimestamp{00:52:55}}
\end{classroomqa}

\section{Четырехвекторы, скаляры и метрика}

Покомпонентные уравнения стали достаточно длинными, чтобы оправдать более удобные обозначения.
Запишем координаты пространства-времени как
\begin{equation}
x^\mu=(x^0,x^1,x^2,x^3)=(t,x,y,z),
\label{eq:sr5-four-coordinate}
\end{equation}
где греческие индексы пробегают все четыре компоненты. Набор из четырех компонент не является четырехвектором автоматически. Он является четырехвектором только в том случае, если его компоненты преобразуются при преобразованиях Лоренца так же, как \(x^\mu\).

\begin{figure}[htbp]
\centering
\includegraphics[width=0.90\textwidth]{lecture_05_figure_15.png}
\caption{Четыре координаты пространства-времени, собранные в \(x^\mu\).}
\lectureframe{00:55:00}
\end{figure}

Утверждение о том, что три пространственные компоненты равны нулю, может быть истинным в одной системе отсчета и ложным в другой. Напротив, если равны нулю все четыре компоненты, все наблюдатели, связанные преобразованиями Лоренца, согласны с тем, что вектор равен нулю. Разности, такие как \(dx^\mu\), тоже являются четырехвекторами.

Инвариантный интервал собственного времени равен
\begin{equation}
d\tau^2
=
dt^2-dx^2-dy^2-dz^2.
\label{eq:sr5-proper-time}
\end{equation}
Он имеет одно значение, а не четыре зависящие от системы отсчета компоненты, поэтому является скаляром.

\begin{figure}[htbp]
\centering
\includegraphics[width=0.90\textwidth]{lecture_05_figure_16.png}
\caption{Четырехсмещение и инвариантный интервал собственного времени.}
\lectureframe{00:57:25}
\end{figure}

\subsection{Метрика Минковского}

Выберем сигнатуру с преобладанием плюсов
\begin{equation}
\eta_{\mu\nu}
=
\operatorname{diag}(-1,1,1,1).
\label{eq:sr5-metric}
\end{equation}
Альтернативная конвенция с преобладанием минусов столь же допустима, но во всех уравнениях необходимо последовательно использовать одну и ту же конвенцию. При выборе \eqref{eq:sr5-metric},
\begin{equation}
dx^\mu dx_\mu
=
-dt^2+dx^2+dy^2+dz^2
=
-d\tau^2.
\end{equation}
Общий знак не меняет инвариантности.

\begin{figure}[htbp]
\centering
\includegraphics[width=0.90\textwidth]{lecture_05_figure_17.png}
\caption{Метрика Минковского с преобладанием плюсов, записанная как диагональная матрица.}
\lectureframe{01:00:30}
\end{figure}

Метрика опускает индекс:
\begin{equation}
A_\mu=\eta_{\mu\nu}A^\nu.
\label{eq:sr5-lower-index}
\end{equation}
Для этой декартовой метрики
\begin{equation}
A^\mu=(A^0,A^1,A^2,A^3),
\qquad
A_\mu=(-A^0,A^1,A^2,A^3).
\label{eq:sr5-upper-lower-components}
\end{equation}
Повторное применение той же метрики поднимает индекс, поскольку \(\eta^{\mu\rho}\eta_{\rho\nu}=\delta^\mu{}_\nu\).

\begin{figure}[htbp]
\centering
\includegraphics[width=0.90\textwidth]{lecture_05_figure_18.png}
\caption{Опускание индекса как умножение на матрицу \(\eta_{\mu\nu}\).}
\lectureframe{01:03:00}
\end{figure}

\begin{classroomqa}
\audiencequestion{Действительно ли здесь обычное умножение матриц содержит суммирование?}
\lecturerresponse{Да. Каждая компонента результата — это строка \(\eta\), умноженная на весь вектор: сумма по столбцовому индексу. Диагональная метрика делает вычисление внешне тривиальным, но для общей матрицы сумма видна явно. \lecturetimestamp{01:04:56}}
\end{classroomqa}

\subsection{Суммирование по Эйнштейну и свободные индексы}

Соглашение Эйнштейна опускает знак суммы всякий раз, когда индекс повторяется один раз вверху и один раз внизу:
\begin{equation}
\eta_{\mu\nu}A^\nu
\equiv
\sum_{\nu=0}^{3}\eta_{\mu\nu}A^\nu.
\end{equation}
Здесь \(\nu\) — немой индекс; по нему выполняется суммирование, и его можно переименовать. Индекс \(\mu\) свободен и обозначает четыре компоненты результата. Обычно индекс не должен появляться более двух раз в одном мономе.

Традиционные названия — \emph{контравариантный} для верхнего индекса и \emph{ковариантный} для нижнего. В лекции намеренно используется более простое обозначение \emph{верхний} и \emph{нижний}, но обе терминологии относятся к одним и тем же правилам преобразования.

\section{Свертки и производные скаляра}

\subsection{Скалярное произведение Минковского}

Свёртка вектора с соответствующей формой с нижним индексом даёт
\begin{equation}
A^\mu A_\mu
=
\eta_{\mu\nu}A^\mu A^\nu
=
-(A^0)^2+(A^1)^2+(A^2)^2+(A^3)^2.
\label{eq:sr5-vector-square}
\end{equation}
Никакой дополнительный минус вручную вставлять не нужно; он уже возникает из временной компоненты с нижним индексом.

\begin{figure}[htbp]
\centering
\includegraphics[width=0.90\textwidth]{lecture_05_figure_19.png}
\caption{Свёртка \(A^\mu A_\mu\), разложенная по временной и пространственным компонентам.}
\lectureframe{01:09:30}
\end{figure}

\begin{classroomqa}
\audiencequestion{Следует ли изображать вектор с нижним индексом как строку, а не как столбец?}
\lecturerresponse{В элементарной матричной записи его можно представлять строкой, но индексная запись не требует такого выбора. Его определяющий признак — то, как он преобразуется и как свёртывается с вектором с верхним индексом. \lecturetimestamp{01:10:26}}
\end{classroomqa}

\begin{classroomqa}
\audiencequestion{Всегда ли греческие индексы относятся ко всем четырём компонентам?}
\lecturerresponse{В этих обозначениях — да. Греческие индексы, такие как \(\mu,\nu\), пробегают по пространству-времени, тогда как латинские индексы, такие как \(i,j\), обычно пробегают только по трём пространственным компонентам. \lecturetimestamp{01:12:23}}
\end{classroomqa}

Та же конструкция работает и для двух различных четырёхвекторов:
\begin{equation}
A^\mu B_\mu
\end{equation}
это скаляр. Одно быстрое доказательство использует тождество поляризации. Поскольку
\((A+B)^2\) и \((A-B)^2\) — обе скалярные величины,
\begin{equation}
(A+B)^2-(A-B)^2=4A^\mu B_\mu
\end{equation}
тоже является скаляром. Это минковская версия обычного скалярного произведения.

\begin{classroomqa}
\audiencequestion{Является ли \(A^\mu B_\mu\) аналогом обычного скалярного произведения?}
\lecturerresponse{Именно, за исключением того, что метрика Минковского задаёт относительный знак между временной и пространственными компонентами. \lecturetimestamp{01:15:44}}
\end{classroomqa}

\subsection{Градиент — ковектор}

Пусть \(\phi(x)\) — скалярное поле, причём \(x\) теперь обозначает все четыре
пространственно-временные координаты. Между соседними событиями
\begin{equation}
d\phi
=
\frac{\partial\phi}{\partial x^\mu}dx^\mu
=
\partial_\mu\phi\,dx^\mu.
\label{eq:sr5-dphi}
\end{equation}
Левая часть — скаляр, а \(dx^\mu\) — четырёхвектор. Следовательно, коэффициент
\begin{equation}
\partial_\mu\phi
=
\left(
\partial_t\phi,
\partial_x\phi,
\partial_y\phi,
\partial_z\phi
\right)
\label{eq:sr5-gradient-lower}
\end{equation}
должен преобразовываться как ковектор. Поднимая его индекс, получаем
\begin{equation}
\partial^\mu\phi
=
\left(
-\partial_t\phi,
\partial_x\phi,
\partial_y\phi,
\partial_z\phi
\right).
\label{eq:sr5-gradient-upper}
\end{equation}

\begin{figure}[htbp]
\centering
\includegraphics[width=0.90\textwidth]{lecture_05_figure_20.png}
\caption{Дифференциал \(d\phi=\partial_\mu\phi\,dx^\mu\), из которого следует ковекторный характер градиента скалярного поля.}
\lectureframe{01:20:00}
\end{figure}

\begin{classroomqa}
\audiencequestion{Превращается ли верхний индекс в \(x^\mu\) в нижний, когда он появляется в знаменателе \(\partial/\partial x^\mu\)?}
\lecturerresponse{Да. Координатный дифференциал \(dx^\mu\) преобразуется контравариантно, поэтому его дуальная производная \(\partial_\mu=\partial/\partial x^\mu\) преобразуется ковариантно. Именно поэтому \(d\phi=\partial_\mu\phi\,dx^\mu\) является скаляром. \lecturetimestamp{01:20:00}}
\end{classroomqa}

\subsection{Инвариантная плотность скалярного поля}

Свернём два вектора производных:
\begin{align}
\partial_\mu\phi\,\partial^\mu\phi
&=
-(\partial_t\phi)^2
+(\partial_x\phi)^2
+(\partial_y\phi)^2
+(\partial_z\phi)^2.
\end{align}
Следовательно, плотность свободного скалярного поля можно компактно записать как
\begin{equation}
\boxed{
\mathcal L_0
=
-\frac12
\partial_\mu\phi\,\partial^\mu\phi
}.
\label{eq:sr5-compact-density}
\end{equation}
Это совпадает с выражением по компонентам в
\eqref{eq:sr5-free-field-action}.

\begin{figure}[htbp]
\centering
\includegraphics[width=0.90\textwidth]{lecture_05_figure_21.png}
\caption{Кинетическая плотность скалярного поля, сведённая к одному свёрнутому выражению.}
\lectureframe{01:25:10}
\end{figure}

\begin{classroomqa}
\audiencequestion{Является ли свёрнутое выражение через первые производные лапласианом?}
\lecturerresponse{Нет. Лапласиан содержит вторые пространственные производные. Это выражение — плотность лагранжиана, построенная из произведений первых производных. \lecturetimestamp{01:24:38}}
\end{classroomqa}

Значение скалярной плотности в физическом событии не зависит от системы отсчёта. Собственное преобразование Лоренца также имеет определитель, единичный по модулю, так что мера \(d^4x\) инвариантна. Следовательно,
\begin{equation}
S=\int d^4x\,\mathcal L
\end{equation}
инвариантно. Независимое от системы отсчёта действие приводит к уравнениям Эйлера--Лагранжа, имеющим один и тот же вид во всех инерциальных системах отсчёта.

\section{Предел Галилея}

\begin{classroomqa}
\audiencequestion{Требовалось ли также, чтобы ньютоновский лагранжиан соблюдал принцип относительности?}
\lecturerresponse{Да, хотя ранее это не подчёркивалось. Ньютоновы уравнения инвариантны относительно преобразований Галилея, предела \(c\to\infty\) преобразований Лоренца. В этом пределе исчезает относительность одновременности. \lecturetimestamp{01:27:16}}
\end{classroomqa}

Чтобы увидеть этот предел, вернём \(c\) в формулы буста вдоль \(x\):
\begin{equation}
x'
=
\frac{x-vt}{\sqrt{1-v^2/c^2}},
\qquad
t'
=
\frac{t-vx/c^2}{\sqrt{1-v^2/c^2}}.
\label{eq:sr5-lorentz-c}
\end{equation}
Переход \(c\to\infty\) даёт
\begin{equation}
\boxed{
x'=x-vt,
\qquad
t'=t
}.
\label{eq:sr5-galilean}
\end{equation}
Это преобразования Галилея. Ньютоновская механика сохраняет принцип относительности, согласно которому законы имеют один и тот же вид во всех инерциальных системах отсчёта,
но при этом она также предполагает абсолютное время, общее для всех наблюдателей.

\begin{figure}[htbp]
\centering
\includegraphics[width=0.90\textwidth]{lecture_05_figure_22.png}
\caption{Лоренцев буст до перехода \(c\to\infty\), после которого остаются галилеевы правила \(x'=x-vt\) и \(t'=t\).}
\lectureframe{01:29:40}
\end{figure}

\begin{classroomqa}
\audiencequestion{Отличается ли преобразование векторов с верхним и нижним индексом только сменой знака?}
\lecturerresponse{Для фиксированной декартовой метрики, используемой здесь, опускание индекса у уже известного вектора меняет только знак его временной компоненты. Однако их законы преобразования Лоренца являются дуальными законами, а не просто произвольным соглашением о знаке; именно метрика переводит одно в другое. \lecturetimestamp{01:31:19}}
\end{classroomqa}

\begin{classroomqa}
\audiencequestion{Является ли всякая с виду обычная функция координат скаляром?}
\lecturerresponse{Нет. Недостаточно того, что в каждой координатной точке задано одно число; это число должно представлять одну и ту же физическую величину в одном и том же событии во всех системах отсчёта. Плотность энергии — контрпример: при лоренцевом бусте ящика меняются и измеренная энергия, и измеренный объём, поэтому наблюдатели не приписывают одну и ту же плотность. \lecturetimestamp{01:33:23}}
\end{classroomqa}

Ещё один пример даёт отдельная компонента четырёхскорости. Сама четырёхскорость является вектором, но её временная компонента меняется при бусте и сама по себе не является скаляром.

\section{Поле Клейна--Гордона}

\subsection{Полевой аналог гармонического осциллятора}

Вернёмся к скалярной плотности и добавим простейший квадратичный потенциал:
\begin{equation}
\boxed{
\mathcal L
=
\frac12
\left[
(\partial_t\phi)^2
-|\boldsymbol{\nabla}\phi|^2
-\mu^2\phi^2
\right]
}.
\label{eq:sr5-kg-density}
\end{equation}
Это полевой аналог
\begin{equation}
L_{\mathrm{osc}}
=
\frac12\dot q^2-\frac12\mu^2q^2.
\end{equation}
Каждая пространственная точка локально ведёт себя как осциллятор, а градиентный член связывает соседние точки.

\begin{figure}[htbp]
\centering
\includegraphics[width=0.90\textwidth]{lecture_05_figure_23.png}
\caption{Квадратичная плотность скалярного поля и её прототип — гармонический осциллятор.}
\lectureframe{01:38:20}
\end{figure}

Уравнение Эйлера--Лагранжа для поля имеет вид
\begin{equation}
\partial_\mu
\left[
\frac{\partial\mathcal L}
{\partial(\partial_\mu\phi)}
\right]
-
\frac{\partial\mathcal L}{\partial\phi}
=0.
\label{eq:sr5-field-el}
\end{equation}
Применяя его к \eqref{eq:sr5-kg-density}, получаем
\begin{equation}
\boxed{
\partial_t^2\phi-\nabla^2\phi+\mu^2\phi=0
}.
\label{eq:sr5-kg}
\end{equation}
Это уравнение Клейна--Гордона. При соглашении о сигнатуре с преимущественно плюсами и определении
\(\Box=\partial_\mu\partial^\mu=-\partial_t^2+\nabla^2\), то же уравнение имеет вид
\((\Box-\mu^2)\phi=0\).

\begin{figure}[htbp]
\centering
\includegraphics[width=0.90\textwidth]{lecture_05_figure_24.png}
\caption{Уравнение Клейна--Гордона, дописанное под плотностью, из которой оно следует.}
\lectureframe{01:40:00}
\end{figure}

Исторически это уравнение сначала исследовалось как релятивистское волновое уравнение для одной частицы. Такая интерпретация столкнулась с трудностями,
особенно в отношении положительной плотности вероятности. Его корректный современный контекст —
релятивистская теория поля, где оно описывает свободные скалярные поля и свободный предел взаимодействующих скалярных полей.

\subsection{Плосковолновые моды}

Поскольку \eqref{eq:sr5-kg} линейно и имеет постоянные коэффициенты, будем искать решение в виде плоской
волны
\begin{equation}
\phi(x)
=
\exp\!\left[
-i\omega t+i\mathbf{k}\cdot\mathbf{x}
\right].
\label{eq:sr5-plane-wave}
\end{equation}
При
\begin{equation}
k_\mu=(-\omega,k_x,k_y,k_z),
\end{equation}
показатель экспоненты равен \(ik_\mu x^\mu\).

\begin{figure}[htbp]
\centering
\includegraphics[width=0.90\textwidth]{lecture_05_figure_25.png}
\caption{Плосковолновой анзац, записанный по компонентам и в виде \(e^{ik_\mu x^\mu}\).}
\lectureframe{01:44:00}
\end{figure}

Дифференцирование даёт
\begin{equation}
\partial_t^2\phi=-\omega^2\phi,
\qquad
\nabla^2\phi=-|\mathbf{k}|^2\phi.
\end{equation}
Подстановка в \eqref{eq:sr5-kg} даёт
\begin{equation}
\left(
-\omega^2+|\mathbf{k}|^2+\mu^2
\right)\phi=0.
\end{equation}
Следовательно, для ненулевой моды требуется
\begin{equation}
\boxed{
\omega^2=|\mathbf{k}|^2+\mu^2
},
\qquad
\omega=\pm\sqrt{|\mathbf{k}|^2+\mu^2}.
\label{eq:sr5-dispersion}
\end{equation}

\begin{figure}[htbp]
\centering
\includegraphics[width=0.90\textwidth]{lecture_05_figure_26.png}
\caption{Подстановка плоской волны и получающееся релятивистское дисперсионное соотношение.}
\lectureframe{01:47:00}
\end{figure}

В естественных единицах \(\mu\) играет роль массы. Восстановление констант для квантового возбуждения,
\begin{equation}
E=\hbar\omega,
\qquad
\mathbf{p}=\hbar\mathbf{k},
\qquad
\mu=\frac{mc}{\hbar},
\end{equation}
превращает \eqref{eq:sr5-dispersion} в
\begin{equation}
E^2=p^2c^2+m^2c^4.
\end{equation}
Ветви с положительной и отрицательной частотами обе входят в пространство классических решений. Общие решения представляют собой суперпозиции этих мод.

\begin{classroomqa}
\audiencequestion{Откуда взялось уравнение Klein--Gordon?}
\lecturerresponse{Оно непосредственно следует из квадратичной лоренц-скалярной плотности через уравнение Эйлера--Лагранжа для поля. Этот пример был выбран потому, что такая простая квадратичная плотность даёт линейное уравнение поля. \lecturetimestamp{01:47:59}}
\end{classroomqa}

\begin{classroomqa}
\audiencequestion{Что в природе подчиняется уравнению Klein--Gordon?}
\lecturerresponse{Свободное классическое скалярное поле. Флуктуации поля Хиггса являются скалярными и в свободном приближении сводятся к форме Klein--Gordon, хотя реальное поле Хиггса также имеет взаимодействия и калибровочные связи. \lecturetimestamp{01:50:41}}
\end{classroomqa}

\section{Скалярный фон как сдвиг массы}

Теперь действие частицы снова входит в рассмотрение:
\begin{equation}
S_{\mathrm p}
=
-\int d\tau\,
\bigl[m+g\phi(x)\bigr].
\label{eq:sr5-final-particle-action}
\end{equation}
Если \(\phi\) имеет постоянное фоновое значение \(\phi_0\), частица движется так,
как если бы её масса была равна
\begin{equation}
\boxed{
m_{\mathrm{eff}}=m+g\phi_0
}.
\label{eq:sr5-effective-mass}
\end{equation}
Если голый член \(m\) обращается в нуль, скалярный фон может полностью задавать
коэффициент при собственном времени.

\begin{figure}[htbp]
\centering
\includegraphics[width=0.90\textwidth]{lecture_05_figure_27.png}
\caption{Зависящее от поля действие частицы, приведённое выше, и компактное уравнение Эйлера--Лагранжа для поля сводят воедино две стороны лекции.}
\lectureframe{01:51:30}
\end{figure}

Это намеренно упрощённая модель хиггсовского порождения массы. В
Стандартной модели массы фермионов возникают из-за связей Юкавы с вакуумным
средним значением поля Хиггса, схематически \(m_f\propto y_f v\), тогда как массы \(W\) и
\(Z\) возникают из калибровочно-ковариантного кинетического члена поля Хиггса. Эта игрушечная
связь на мировой линии передаёт основную идею--ненулевой скалярный фон
изменяет инерциальную массу--не воспроизводя полной электрослабой теории.
Установленная ранее взаимность остаётся в силе: варьирование того же
взаимодействия по полю делает частицу источником. Поле действует на
частицы, а частицы в свою очередь действуют на поле.

\subsection{Почему комплексные экспоненты полезны}

Поле не обязано быть физически комплексным. Формула Эйлера даёт
\begin{equation}
e^{i(\omega t-\mathbf{k}\cdot\mathbf{x})}
=
\cos(\omega t-\mathbf{k}\cdot\mathbf{x})
+i\sin(\omega t-\mathbf{k}\cdot\mathbf{x}).
\end{equation}
Поскольку уравнение вещественно и линейно, его действительная и мнимая части
по отдельности являются вещественными решениями. Это волны, сдвинутые на
четверть цикла, и любая вещественная линейная комбинация снова является решением.

\begin{classroomqa}
\audiencequestion{Когда берут действительную часть, не складывается ли мнимая часть каким-то образом в нуль?}
\lecturerresponse{Нет. Комплексная экспонента объединяет два независимых вещественных решения, косинус и синус. Можно сохранить любую из частей или образовать вещественные комбинации. Экспонента используется потому, что дифференцирование превращается в умножение. \lecturetimestamp{01:53:29}}
\end{classroomqa}

\begin{figure}[htbp]
\centering
\includegraphics[width=0.90\textwidth]{lecture_05_figure_28.png}
\caption{Комплексная плоская волна, разложенная на косинус и синус и изображённая как две вещественные волны, сдвинутые по фазе.}
\lectureframe{01:56:00}
\end{figure}

\begin{classroomqa}
\audiencequestion{Является ли масса, приобретаемая электромагнитным полем в сверхпроводнике, явлением того же рода?}
\lecturerresponse{Это тесно связано: в сверхпроводнике конденсат сообщает электромагнитному полю эффективную массу в среде, что приводит к эффекту Мейсснера. Полное объяснение требует рассмотрения электромагнитного поля и откладывается. \lecturetimestamp{01:56:14}}
\end{classroomqa}

\begin{classroomqa}
\audiencequestion{Можно ли представлять себе механизм Хиггса как частицу Хиггса, летящую через пространство и раздающую массу?}
\lecturerresponse{Нет. Существенным объектом является фоновое значение поля Хиггса, а не поток отдельных частиц Хиггса. В игрушечном действии фон просто изменяет коэффициент, который играет роль массы частицы. \lecturetimestamp{01:57:31}}
\end{classroomqa}

\begin{classroomqa}
\audiencequestion{Могут ли разные связи с одним и тем же скалярным фоном порождать разные массы?}
\lecturerresponse{Да. Сдвиг пропорционален константе связи \(g\). В Стандартной модели именно разные связи Юкавы и объясняют, почему разные фермионы приобретают разные массы от одного и того же фона Хиггса. \lecturetimestamp{01:58:04}}
\end{classroomqa}

\begin{classroomqa}
\audiencequestion{Является ли \(\phi\) само по себе массовым членом?}
\lecturerresponse{В этой игрушечной модели \(g\phi_0\) во всех механических отношениях функционирует как вклад в коэффициент массы. Необходимо учитывать и константу связи, и фоновое значение. \lecturetimestamp{01:58:42}}
\end{classroomqa}

\begin{classroomqa}
\audiencequestion{Если бы фон Хиггса исчез, стала бы каждая частица безмассовой?}
\lecturerresponse{Нет. Электроны, кварки и слабые калибровочные бозоны потеряли бы свои обычные электрослабые вклады в массу, причём детали для нейтрино зависят от механизма, который сообщает нейтрино массу. Сам скалярный сектор не обязан становиться безмассовым лишь потому, что его фоновое значение обращается в нуль. Большая часть массы протона и нейтрона обусловлена вместо этого КХД-связыванием и конфайнментом, так что она не генерируется непосредственно полем Хиггса. \lecturetimestamp{01:58:51}}
\end{classroomqa}

Тем не менее следствие для атомов радикально. В естественных единицах боровский масштаб ведёт себя как
\begin{equation}
a_0\sim\frac{1}{\alpha m_e}
\end{equation}
При стремлении массы электрона к нулю обычный атомный масштаб длины расходится; обычные локализованные атомы перестали бы существовать.
Протон и нейтрон были бы затронуты гораздо менее прямо, потому что их
масса является преимущественно динамической энергией КХД.

Теперь все части складываются в единую связную линию рассуждения. Геометрия
Лоренца определяет наблюдаемые длины волн. Один смешанный член в одном действии заставляет
две динамические системы влиять друг на друга. Скалярное поле становится
одновременно потенциалом для частицы и полем, источником которого служит эта частица. Метрика
и индексные обозначения делают лоренц-инвариантность наглядной, а простейшая
квадратичная инвариантная плотность порождает волны Klein--Gordon. Наконец, постоянный
скалярный фон превращает ту же самую связь в сдвиг массы. И на всём протяжении
сохраняется один и тот же метод: построить инвариантное действие, проварьировать
каждый динамический объект и прочитать взаимную физику из получающихся уравнений.